% \iffalse meta-comment
%
% hit-keylist.dtx 是各个类的键清单：这个类接受哪些信息字段与常量
%
% \fi
%
% \section{键清单（keylist）}
%
% \changes{v3.2a}{2026/09/13}{元信息字段的声明单独成文件，两个类共用一份
%   \texttt{shared-keylist}：它说的是「有哪些字段」，字段怎么存、旧名怎么兼容那套
%   机制在 \file{src/engine/hit-key-engine.dtx}。原先 96 个名字在两个文件里各写一遍，
%   加一个公共字段要改两处，漏一处没有任何东西会说}
% \changes{v3.2a}{2026/08/16}{\cs{hitsetup} 改用 l3keys 做入口，配置项按职责分族：
%   \texttt{info}（\option{specialty} 改名 \option{major}，加 \option{student-number}
%   别名）、\texttt{term}（\file{cfg} 里的可本地化字符串，加编号之后的间距
%   \option{after-}\meta{层级}\option{-number} 一批）、\texttt{bib}（按
%   kebab-biblatex-gb7714 的词汇重起名，\option{style} 进 \option{bibtex} 子组）、
%   \texttt{style}、\texttt{toc}（挪到共享层，报告也认）、\texttt{math}。
%   旧名照常生效，用一次提示一次；\option{tocblank} 废弃不再起作用}
% \changes{v3.2a}{2026/09/15}{加 \texttt{layout} 族：收 Word「页面设置」与「文档网格」
%   两个对话框上的键，按四个选项卡分组嵌套，转给 \cs{docxlike\_page\_setup:n}；
%   页底策略 \option{bottom} 取 \texttt{ragged}／\texttt{flush}；带部件表，各部件把
%   与全篇的差异写在自己名下；页码线、页眉线、右开页从 \option{style} 挪进来按部件
%   寻址；\option{header}／\option{footer} 的 \option{shift} 与 \option{shown} 照
%   Word 的对话框设；\option{grid}/\option{chars}/\option{latin} 与
%   \option{grid}/\option{enabled}；旧版面下不推部件版面（博士后摘要页曾被推成零边距）}
% \changes{v3.2a}{2026/09/21}{加 \option{style}/\option{em-dash}、
%   \option{style}/\option{fake-slant-songti}、\option{style}/\option{float-sep}、
%   \option{style}/\option{verb-space}}
% \changes{v3.2a}{2026/08/25}{加 \texttt{defense} 族（评阅人、答辩委员会名单与决议正文）、
%   \option{abbr} 族（缩略语词条）、\option{cover/title-en-xiaoer} 三值键（封面与内封
%   同名同义）}
% \changes{v3.2a}{2026/08/29}{日期字段写到哪级排到哪级，封面日期的回退改看存原值
%   那个宏；四个摘要环境改用内核的 \texttt{+b} 参数类型，\pkg{environ} 删掉；
%   \pkg{etoolbox} 只在共享层装一次}
%
% 这个文件只有 \texttt{shared-keylist} 一个守卫。两个类合成一个之后没有「各取各的」
% 这回事了，终稿与报告都从这一份里取，各档要不要某个键由值那一层的条件分，
% 不由守卫分。
%
% 这一层只声明，不给值。值在 \file{src/config/} 那几份，机制在
% \file{hit-key-engine.dtx}，规矩见 CONTRIBUTING §4.2。
%
% ^^A ======================================================================
% ^^A Module shared-keylist: 信息字段与常量的声明
% ^^A ======================================================================
%    \begin{macrocode}
%<*shared-keylist>
\hit@define@title{zh}
%    \end{macrocode}
%
% 两档都有的信息字段：封面上要填的那批，学院、专业、导师、作者、学号。
%
% 专业那个字段 v3.1e 叫 \option{csubject}，注释写的是「中文专业」，v3.2a 起叫
% \option{major}。学科是另一个字段 \option{discipline}（v3.1e 的 \option{cxueke}），
% 封面上「工学博士」的「工学」取的是它，两者不是一回事。
%
% 学号两个名字都对：\option{student-id} 与 \option{student-number}，写哪个都不提示。
%    \begin{macrocode}
\hit@define@term{author-zh}
\hit@define@term{supervisor-zh}
\hit@define@term{affiliation-zh}
\hit@define@term{major-zh}
\hit@define@date@term{date-zh}{zh}
\hit@define@term{student-id}
\hit@define@term@alias{student-number}{student-id}
\hit@declare@info@field
  {
    title-zh , author-zh , supervisor-zh , affiliation-zh , major-zh , date-zh ,
    student-id , student-number
  }
%    \end{macrocode}
%
% 两档都有的常量。三种存法各一批：\cs{hit@declare@constant} 存进
% \cs{hit@}\meta{名}，\cs{hit@declare@constant@as} 驱动同名的既有命令，
% \cs{hit@declare@constant@ctex} 转发给 \cs{ctexset}。后两批两档逐字相同。
%    \begin{macrocode}
\hit@declare@constant@as
  {
    equationname , listequationname
  }
\hit@declare@constant
  {
    abstract-heading-zh , abstract-heading-en ,
    abstract-in-english-toc-zh , abstract-in-english-toc-en ,
    list-of-figures-in-english-toc , list-of-tables-in-english-toc ,
    list-of-equations-in-english-toc
  }
\hit@declare@constant@ctex
  {
    chapter-name = chapter/name , appendixname , contentsname , listfigurename ,
    listtablename , figurename , tablename , bibname , indexname
  }
\hit@declare@constant
  {
    caption-algorithm-zh , caption-algorithm-en , base-campus-full , declarations-authorization ,
    declarations-author-signature , titlepage-affiliation-bachelor-zh , titlepage-date-bachelor-zh , cover-report-date-bachelor , cover-check-date-bachelor ,
    titlepage-student-bachelor , cover-student-bachelor , base-document-type-bachelor , base-document-type-level-bachelor ,
    cover-teacher-comment-bachelor , cover-teacher-signature-bachelor ,
    bibliography-heading-zh , titlepage-affiliation-zh , titlepage-associate-supervisor-zh , titlepage-author-zh ,
    titlepage-co-supervisor-zh , titlepage-date-zh , base-degree-of-discipline-zh ,
    titlepage-degree-applied-zh , base-university-zh , titlepage-institution-zh , titlepage-major-zh , titlepage-supervisor-zh , base-document-type , cover-title-prefix , base-degree-level-zh ,
    declarations-date-blank , mainmatter-eqdenote-label-zh , mainmatter-eqdenote-label-en , mainmatter-eqdenote-dash , caption-figure-prefix-en ,
    caption-table-prefix-en , declarations-date , titlepage-affiliation , cover-affiliation ,
    titlepage-affiliation-major , cover-affiliation-major , titlepage-major , cover-major ,
    titlepage-date , cover-report-date , cover-enroll-date ,
    titlepage-student-id , cover-student-id , titlepage-student , cover-student ,
    titlepage-supervisor , cover-supervisor , cover-title-label ,
    cover-school-bottom-mark-bachelor , cover-school-bottom-mark-harbin-master , cover-school-bottom-mark-harbin-doctor , base-campus , base-hi ,
    titlepage-classified-index-international-zh ,
    titlepage-classified-index-national-zh , base-school-code , titlepage-school-code ,
    base-secrecy , cover-note-interim-shenzhen-doctor ,
    titlepage-major-bachelor-zh , cover-major-bachelor , titlepage-student-id-bachelor , cover-student-id-bachelor ,
    titlepage-supervisor-bachelor-zh , cover-supervisor-bachelor , base-stage-document-type , base-stage-interim , base-stage-proposal ,
    declarations-supervisor-signature , base-title-separator-zh ,
    titlepage-degree-sentence , titlepage-degree-line ,
    titlepage-mark-dissertation , titlepage-mark-practice , cover-bottommark
  }
%    \end{macrocode}
%
% 编号之后排什么。标题四级各一个，图题表题各一个，中英各一份。取值是排版内容，
% 写 \cs{quad}、\cs{enspace} 这种间距最常见，写别的也行。
%
% 默认一律留空，空就退回各处原有的取值，见 \cs{hit@after@number:nnn}。
% \option{after-chapter-number} 只有终稿用得上。报告的底类也是
% \pkg{ctexbook}，\cs{chapter} 在，但版式里不排章。
%    \begin{macrocode}
\hit@declare@constant
  {
    after-chapter-number-zh , after-chapter-number-en ,
    after-section-number-zh , after-section-number-en ,
    after-subsection-number-zh , after-subsection-number-en ,
    after-subsubsection-number-zh , after-subsubsection-number-en ,
    after-figure-caption-number-zh , after-figure-caption-number-en ,
    after-table-caption-number-zh , after-table-caption-number-en
  }
%    \end{macrocode}
%
% 四级标题一起设的写法。写了它等于把四个逐一写一遍，所以后面再单设某一级仍然管用。
%    \begin{macrocode}
\clist_map_inline:nn { zh , en }
  {
    \keys_define:nn { hit / term }
      {
        after-heading-number-#1 .meta:n =
          {
            after-chapter-number-#1       = {##1} ,
            after-section-number-#1       = {##1} ,
            after-subsection-number-#1    = {##1} ,
            after-subsubsection-number-#1 = {##1} ,
          } ,
      }
  }
%    \end{macrocode}
%
% \option{bib}/\option{bibtex}/\option{style} 是 \cs{bibliographystyle} 的参数。
% 它原先混在 \option{struct} 的部件表里，可部件说的是「排不排这一块」，样式
% 说的是「排成什么样」，不是一回事。后来又从 \option{bib} 的根层挪进
% \option{bibtex} 子组：挑 \file{.bst} 只在传统一线有意义，biblatex 一线的
% 样式锁死国标（见 \file{hit-bibstyle.dtx}）。键在共享层声明，默认值在这里给。
%
% \texttt{toc} 族。三个键两档都有，所以声明在共享层：报告一样有目录，一样要
% 排到第几级、西文用不用无衬线体。
%
% 变量在这里声明，早于 \cs{ProcessKeyOptions}，所以旧名的类选项
% 处理器直接写它们就行。取值的含义两档一致，怎么落到版面上各自决定——
% \option{depth} 数的是级数，学位论文有章、报告没有，所以两边换算到
% \cs{tocdepth} 差一。
%
% \option{english} 只有终稿做得了事，报告一份文档只排一份目录，
% 语言跟着类选项 \option{lang} 走。报告在自己那半把这个键改成一条提示。
%    \begin{macrocode}
\tl_new:N \g__hit_bib_style_tl    \tl_gset:Nn \g__hit_bib_style_tl   { hithesis }
\tl_new:N \g__hit_toc_english_tl  \tl_gset:Nn \g__hit_toc_english_tl { auto }
\tl_new:N \g__hit_toc_latin_natural_tl
\tl_gset:Nn \g__hit_toc_latin_natural_tl { auto }
\int_new:N \g__hit_toc_depth_int  \int_gset:Nn \g__hit_toc_depth_int { 3 }
\bool_new:N \g__hit_toc_latin_sans_bool
\keys_define:nn { hit / toc }
  {
    english .code:n = { \__hit_tristate_set:Nn \g__hit_toc_english_tl {#1} } ,
    latin-natural .code:n =
      { \__hit_tristate_set:Nn \g__hit_toc_latin_natural_tl {#1} } ,
    latin-natural .default:n = true ,
    depth   .code:n = { \int_gset:Nn \g__hit_toc_depth_int {#1} } ,
  }
\__hit_boolean_key:nnn { toc } { latin-sans } { toc_latin_sans }
%    \end{macrocode}
%
% 目录条目里的西文字体。开关关着时它展开成空，所以可以无条件套在字体那一串
% 前面，不必在每个用到的地方再判断一次。排版时才读标志位，导言区改得动。
%    \begin{macrocode}
\cs_new:Npn \hit@toc@font
  { \bool_if:NT \g__hit_toc_latin_sans_bool { \sffamily } }
%    \end{macrocode}
%
% 目录条目的四个 \option{styles} 目标，与四级条目一一对应。
% \cs{hit@toc@target:n}\marg{第几级}从 1 数起；终稿的目录级从 0 数（章是 0），
% 报告的正文错开一格（节即章），两边各自加减一次再传进来。
%
% \cs{hit@toc@latin@pick:n}\marg{条目内容}按内容有没有汉字立
% \cs{l\_docxlike\_format\_latin\_pick\_bool}，也就是告诉 \cs{docxlike\_format\_compute:n}
% 这一条的行高按西文字体算。\option{latin-natural} 关掉就一律不立，纯西文的
% 条目也按目标写的那副字算。
%
% \cs{hit@toc@line@metrics:nn}\marg{第几级}\marg{条目内容}是两者合起来那一步：
% 排每一条之前调一次，设好这一条的 \cs{baselineskip}，模型见
% \file{src/docxlike/docxlike-format.dtx} 的 \cs{docxlike\_format\_stack\_line:n}。
%    \begin{macrocode}
\clist_const:Nn \c__hit_toc_target_clist
  { toc-chapter , toc-section , toc-subsection , toc-subsubsection }
\cs_new:Npn \hit@toc@target:n #1
  { \clist_item:Nn \c__hit_toc_target_clist { \int_eval:n {#1} } }
\cs_new_protected:Npn \hit@toc@latin@pick:n #1
  {
    \bool_set_false:N \l_docxlike_format_latin_pick_bool
    \docxlike_format_cjk:nF {#1}
      {
        \__hit_tristate_if:NnT \g__hit_toc_latin_natural_tl { \c_true_bool }
          { \bool_set_true:N \l_docxlike_format_latin_pick_bool }
      }
  }
\cs_new_protected:Npn \hit@toc@line@metrics:nn #1#2
  {
    \hit@toc@latin@pick:n {#2}
    \exp_args:Nx \docxlike_format_stack_line:n { \hit@toc@target:n {#1} }
  }
%    \end{macrocode}
%
% \texttt{layout} 族收 Word 的页面设置。第一个键是 \option{bottom}，取值
% \texttt{ragged}（页面排到哪算哪）与 \texttt{flush}（页底拉齐），名字取自
% \hologo{LaTeX} 自己的两个命令。默认写 \texttt{auto}，眼下等于
% \texttt{ragged}——模板里所有键的默认值都是 \texttt{auto}，这样「没指定」
% 和「指定了某个值」在代码里分得开。
%
% 这一个键同时选定了\emph{按规范的哪一栏排}，因为两者一一对应：
% \begin{longtblr}[caption={\option{bottom} 两档的依据}, label={tab:impl-bottom}]
%   {colspec={l l l}}
% \toprule
% & 依据 & 间距 \\ \midrule
% \texttt{ragged} & 写作指南 3.2 节的行数与多倍值 & 刚性 \\
% \texttt{flush} & 3.8 节的毫米区间 & 区间即伸缩范围 \\
% \bottomrule
% \end{longtblr}
% 3.2 节给的是确切值，刚性；3.8 节给的全是区间（正文 3\,--\,4\,mm、条标题
% 6\,--\,7\,mm 之类），区间本身就是伸缩范围。按 3.2 节排就没有可拉伸的东西，
% 页底拉不齐；按 3.8 节排就有余地。所以不是两个独立的开关，是一件事的两面。
%
% 反过来说，「按毫米排但保持刚性」这一档不存在——毫米栏从来没给过确切值，
% 硬要刚性就得从区间里挑一个数，那个数没有依据。
%
% 正文的余量算得出来：自然值 3.03\,mm、上限 4\,mm，每行富余 0.97\,mm，
% 一页 33 行合计 32\,mm，而页底缺口通常不到 7\,mm。所以实际每行只拉 0.2\,mm
% 上下，肉眼看不出来。
% \emph{\cs{\_\_hit\_layout\_bottom:n}。}
% \option{layout} 是新版面的接口，不该管旧版面：旧版面进了兼容期就锁死，
% 章节格式在那边设了不管用，页底策略在那边设了却管用，说不通。所以这个键
% 只记下取值，能不能落到标志位上由 \cs{\_\_hit\_layout\_bottom\_apply:} 决定。
%
% 类选项阶段 \cs{g\_\_hit\_new\_layout\_bool} 还没算出来（它要等所有类选项读完
% 才知道 \option{v2-layout} 到底开没开），所以那时按兵不动，等选项后处理
% 那里补一次。用 \cs{hitsetup} 设的时候标志位已经有值，当场就能落。
% \file{hit-report.cls} 没有新旧版面之分，那边这个标志位不存在，直接落。
%
% 旧版面要改页底策略走 \option{v2-layout}\texttt{=\{bottom=\ldots\}}。
%    \begin{macrocode}
\tl_new:N \g__hit_layout_bottom_tl
\cs_new_protected:Npn \__hit_layout_bottom:n #1
  {
    \tl_gset:Nn \g__hit_layout_bottom_tl {#1}
    \__hit_layout_bottom_apply:
  }
\cs_new_protected:Npn \__hit_layout_bottom_apply:
  {
    \tl_if_empty:NF \g__hit_layout_bottom_tl
      { \__hit_layout_bottom_check: }
  }
\cs_new_protected:Npn \__hit_layout_bottom_check:
  {
    \bool_if_exist:NTF \g__hit_layout_two_bool
      {
        \bool_if:NF \g__hit_layout_two_bool
          { \__hit_layout_bottom_set: }
      }
      { \__hit_layout_bottom_set: }
  }
\cs_new_protected:Npn \__hit_layout_bottom_set:
  {
    \str_if_eq:VnTF \g__hit_layout_bottom_tl { ragged }
      { \bool_gset_true:N  \g__hit_raggedbottom_bool }
      { \bool_gset_false:N \g__hit_raggedbottom_bool }
    \hit@docxlike@sync:
  }
%    \end{macrocode}
%
%    \begin{macrocode}
\keys_define:nn { hit / layout }
  {
    bottom .choice: ,
    bottom / auto   .code:n = { \__hit_layout_bottom:n { ragged } } ,
    bottom / ragged .code:n = { \__hit_layout_bottom:n { ragged } } ,
    bottom / flush  .code:n = { \__hit_layout_bottom:n { flush  } } ,
    bottom .value_required:n = true ,
  }
%    \end{macrocode}
%
% 其余的键就是 Word 那两个对话框上的格子，一格一个键，换算规则见
% \file{src/docxlike/docxlike-page.dtx}：
% \begin{latex}
% \hitsetup{ layout = {
%   margin = { top = 107.75bp , bottom = 85.05bp , left = 85.05bp , right = 85.05bp } ,
%   header = { from-edge = 85.05bp } , footer = { from-edge = 65.2bp } ,
%   grid = { lines = { pitch = 19.55bp } , chars = { pitch = 12.675bp } } ,
% } }
% \end{latex}
%
% 这些键不当场生效，先攒进一张表，一次 \cs{hitsetup} 里的都收齐了再交给
% \cs{docxlike\_page\_setup:n} 重算一遍。版心的每个数都要另外几个数才算得出来——
% 版心高要上下边距，页眉位置要上边距减页眉距边界——一格一格地即算即用会拿
% 半张表去算，还会因为暂时没有行距而误报。
%
% 表是全局只增的，所以分几次写等同于一次写全，后写的盖先写的。
%
% \option{header}/\option{rule} 那两个键与 \option{footer}/\option{shift} 是把类自己的
% 页眉页脚接到 Word 位置上的，来路见 \file{src/docxlike/docxlike-page.dtx}。
%    \begin{macrocode}
\tl_new:N \g__hit_layout_msword_tl
\tl_new:N \l__hit_layout_target_tl
\tl_set:Nn \l__hit_layout_target_tl { g__hit_layout_msword_tl }
\bool_new:N \l__hit_layout_global_bool
\bool_set_true:N \l__hit_layout_global_bool
\bool_new:N \l__hit_layout_dirty_bool
\cs_new_protected:Npn \__hit_layout_msword:nn #1#2
  {
    \tl_gput_right:cx { \l__hit_layout_target_tl } { #1 = \exp_not:n {#2} , }
    \bool_if:NT \l__hit_layout_global_bool
      { \bool_set_true:N \l__hit_layout_dirty_bool }
  }
%    \end{macrocode}
%
% \cs{\_\_hit\_layout\_leaf:nnn}\marg{块}\marg{键}\marg{docxlike 键} 把一个叶子键
% 定义两遍：一遍在块的子族里，一遍在 \texttt{hit\,/\,layout} 里带全路径。所以
% 下面两种写法等价，嵌套着写和顺手写一条都行：
% \begin{latex}
% layout = { header = { from-edge = 85.05bp } }
% layout = { header/from-edge = 85.05bp }
% \end{latex}
% \meta{docxlike 键}写成 \pkg{docxlike} 页面设置族里的路径，如 \texttt{margins\,/\,top}，
% 攒进表里是 \texttt{margins/top = 85.05bp} 这样一条。
%    \begin{macrocode}
\cs_new_protected:Npn \__hit_layout_leaf:nnn #1#2#3
  {
    \keys_define:nn { hit / layout / #1 }
      {
        #2 .code:n = { \__hit_layout_msword:nn {#3} {##1} } ,
        #2 .value_required:n = true ,
      }
    \keys_define:nn { hit / layout }
      {
        #1 / #2 .code:n = { \__hit_layout_msword:nn {#3} {##1} } ,
        #1 / #2 .value_required:n = true ,
      }
  }
\cs_new_protected:Npn \__hit_layout_block:nn #1#2
  {
    \keys_define:nn { hit / layout / #1 }
      {
        #2 .code:n = { \hit@keys@set:nn { hit / layout / #1 / #2 } {##1} } ,
        #2 .value_required:n = true ,
      }
    \keys_define:nn { hit / layout }
      {
        #1 / #2 .code:n = { \hit@keys@set:nn { hit / layout / #1 / #2 } {##1} } ,
        #1 / #2 .value_required:n = true ,
      }
  }
\clist_map_inline:nn { paper , margin , header , footer , grid }
  {
    \keys_define:nn { hit / layout }
      {
        #1 .code:n = { \hit@keys@set:nn { hit / layout / #1 } {##1} } ,
        #1 .value_required:n = true ,
      }
  }
%    \end{macrocode}
%
% \emph{按部件覆盖。} Word 的一个分节对应我们的一个部件：封面与内封的
% \texttt{docGrid} 行跨度是 395 缇，摘要目录 391，正文 383。写在部件名下的键
% 只对那个部件生效，没写的继承全篇。
%
% 全篇的取值攒在 \cs{g\_\_hit\_layout\_msword\_tl}，各部件攒在
% \cs{g\_\_hit\_layout\_msword\_}\meta{部件}\cs{\_tl}。收键时把
% \cs{l\_\_hit\_layout\_target\_tl} 指向哪个，就攒进哪个。部件那一路不置
% \emph{dirty}——它不即时生效，等排到那个部件再推。
%    \begin{macrocode}
\clist_const:Nn \c__hit_layout_part_clist
  { cover , titlepage , frontmatter , mainmatter , backmatter }
\clist_map_inline:Nn \c__hit_layout_part_clist
  {
    \tl_new:c { g__hit_layout_msword_#1_tl }
    \tl_new:c { g__hit_openright_#1_tl }
    \tl_gset:cn { g__hit_openright_#1_tl } { auto }
    \tl_new:c { g__hit_openright_default_#1_tl }
    \tl_gset:cn { g__hit_openright_default_#1_tl } { false }
  }
\tl_new:N \g__hit_openright_tl
\tl_gset:Nn \g__hit_openright_tl { auto }
\tl_new:N \l__hit_layout_part_tl
\keys_define:nn { hit / layout }
  {
    openright .code:n =
      {
        \tl_if_empty:NTF \l__hit_layout_part_tl
          { \tl_gset:Nn \g__hit_openright_tl {#1} }
          {
            \tl_gset:cn { g__hit_openright_ \l__hit_layout_part_tl _tl } {#1}
          }
      } ,
    openright .value_required:n = true ,
%    \end{macrocode}
%
% \option{openright-default} 只给默认值表用。它写的是最底下那一层：用户在部件
% 或全篇写了值，都盖得住它。学位耦合就住在这一层，见
% \file{src/config/} 下那几份。
%    \begin{macrocode}
    openright-default .code:n =
      {
        \tl_if_empty:NF \l__hit_layout_part_tl
          {
            \tl_gset:cn
              { g__hit_openright_default_ \l__hit_layout_part_tl _tl } {#1}
          }
      } ,
    openright-default .value_required:n = true ,
%    \end{macrocode}
%
% \option{page-number-line} 决定那一段的页码要不要单独占一行。它本来就在按段
% 分，挪进来之后与 \option{openright} 共用部件名，不必另起一套写法。
%    \begin{macrocode}
    page-number-line .code:n =
      {
        \tl_if_empty:NF \l__hit_layout_part_tl
          {
            \cs_if_exist:cT
              { g__hit_page_number_line_ \l__hit_layout_part_tl _tl }
              {
                \__hit_tristate_gset:cn
                  { g__hit_page_number_line_ \l__hit_layout_part_tl _tl } {#1}
              }
          }
      } ,
    page-number-line .value_required:n = true ,
  }
\clist_map_inline:Nn \c__hit_layout_part_clist
  {
    \keys_define:nn { hit / layout }
      {
        #1 .code:n =
          {
            \group_begin:
              \tl_set:Nn \l__hit_layout_target_tl { g__hit_layout_msword_#1_tl }
              \tl_set:Nn \l__hit_layout_part_tl { #1 }
              \bool_set_false:N \l__hit_layout_global_bool
              \hit@keys@set:nn { hit / layout } {##1}
            \group_end:
          } ,
        #1 .value_required:n = true ,
      }
  }
%    \end{macrocode}
%
% \cs{\_\_hit\_layout\_msword\_active:TF} 这份文档走不走 Word 那套版面。旧版面的尺寸写在
% \file{src/flow/hit-flow-geometry.dtx} 里，\option{layout} 那一族的默认值整段不发，
% 攒出来的键值表是空的，拿空表去算只会得到一页零边距。报告没有这个标志位，
% 那边只有一套版面，当作走。
%    \begin{macrocode}
\prg_new_protected_conditional:Npnn \__hit_layout_msword_active: { T , TF }
  {
    \bool_if_exist:NTF \g__hit_layout_two_bool
      {
        \bool_if:NTF \g__hit_layout_two_bool
          { \prg_return_false: } { \prg_return_true: }
      }
      { \prg_return_true: }
  }
%    \end{macrocode}
%
% \begin{macro}{\hit@layout@push@part:n}
% \begin{macro}{\hit@layout@pop@part:}
% 排某个部件之前换上它那一节的版面，排完 \cs{hit@layout@pop@part:} 还回去。
% 传进去的是「全篇取值 + 该部件的差异」，后写的覆盖先写的。
%
% 两条都跟着 \cs{\_\_hit\_layout\_msword\_active:TF} 走，一起做或者一起不做，
% 不然 \cs{docxlike\_page\_setup\_pop:} 会拿上一次存的值往回还。
%    \begin{macrocode}
\cs_new_protected:Npn \hit@layout@push@part:n #1
  {
    \__hit_layout_msword_active:T
      {
        \use:x
          {
            \exp_not:N \docxlike_page_setup_push:n
              {
                \exp_not:v { g__hit_layout_msword_tl }
                \exp_not:v { g__hit_layout_msword_#1_tl }
              }
          }
      }
  }
\cs_new_protected:Npn \hit@layout@pop@part:
  { \__hit_layout_msword_active:T { \docxlike_page_setup_pop: } }
%    \end{macrocode}
% \end{macro}
% \end{macro}
%    \begin{macrocode}
\__hit_layout_leaf:nnn { paper } { width }  { paper / width }
\__hit_layout_leaf:nnn { paper } { height } { paper / height }
\__hit_layout_leaf:nnn { margin } { top }    { margins / top }
\__hit_layout_leaf:nnn { margin } { bottom } { margins / bottom }
\__hit_layout_leaf:nnn { margin } { left }   { margins / left }
\__hit_layout_leaf:nnn { margin } { right }  { margins / right }
\__hit_layout_leaf:nnn { header } { from-edge } { layout / from-edge / header }
\__hit_layout_leaf:nnn { footer } { from-edge } { layout / from-edge / footer }
\__hit_layout_block:nn { header } { rule }
\__hit_layout_block:nn { footer } { shift }
\__hit_layout_block:nn { grid } { chars }
\__hit_layout_block:nn { grid } { lines }
%    \end{macrocode}
%
% \option{header}/\option{rule} 收两种写法：给 \texttt{true} / \texttt{false} /
% \texttt{auto} 是开关，给键值是调参数。仓库里 \option{open-right} 与
% \option{student-type} 都是这个套路，判有没有等号。
%
% \begin{quote}
% \cs{hitsetup}\{ layout = \{ header = \{ rule = false \} \} \}\\
% \cs{hitsetup}\{ layout = \{ header = \{ rule = \{ width = 2.25bp \} \} \} \}
% \end{quote}
%    \begin{macrocode}
\keys_define:nn { hit / layout / header }
  {
    rule .code:n =
      {
        \tl_if_in:nnTF {#1} { = }
          { \hit@keys@set:nn { hit / layout / header / rule } {#1} }
          { \__hit_tristate_set:Nn \g__hit_header_rule_tl {#1} }
      } ,
    rule .value_required:n = true ,
  }
\keys_define:nn { hit / layout }
  {
    header / rule .code:n =
      {
        \tl_if_in:nnTF {#1} { = }
          { \hit@keys@set:nn { hit / layout / header / rule } {#1} }
          { \__hit_tristate_set:Nn \g__hit_header_rule_tl {#1} }
      } ,
    header / rule .value_required:n = true ,
  }
\__hit_layout_leaf:nnn { header / rule } { width }     { header-rule-width }
\__hit_layout_leaf:nnn { header / rule } { from-text } { header-rule-from-text }
\__hit_layout_leaf:nnn { footer / shift } { vertical }   { footer-shift-vertical }
\__hit_layout_leaf:nnn { footer / shift } { horizontal } { footer-shift-horizontal }
%    \end{macrocode}
%
% 页眉页脚各一条记录，与 Word「页面设置」和页眉段落的对话框同形：
% \option{from-edge} 距边界、\option{shown} 排不排（\texttt{auto}/\texttt{true}/\texttt{false}）、
% \option{rule} 页眉段落的下边框（\option{width}、\option{from-text}，或直接给开关）、
% \option{shift} 整块挪动（\option{vertical} 往下为正、\option{horizontal} 往右为正，
% 是量出来的补偿，不折缇）。字体字号行距在 \option{styles}/\option{Header} 与
% \option{styles}/\option{Footer} 上，页眉那一行的行盒按它算。部件可以按名覆盖：
% \begin{quote}
% \cs{hitsetup}\{ layout = \{ header = \{ shown = false \}, frontmatter = \{ header = \{ shown = true \} \} \} \}
% \end{quote}
%    \begin{macrocode}
\__hit_layout_block:nn { header } { shift }
\__hit_layout_leaf:nnn { header / shift } { vertical }   { header-shift-vertical }
\__hit_layout_leaf:nnn { header / shift } { horizontal } { header-shift-horizontal }
\__hit_layout_leaf:nnn { header } { shown } { header-shown }
\__hit_layout_leaf:nnn { footer } { shown } { footer-shown }
\__hit_layout_leaf:nnn { grid / chars } { pitch }    { document-grid / characters / pitch }
\__hit_layout_leaf:nnn { grid / chars } { per-line } { document-grid / characters / characters-per-line }
\__hit_layout_leaf:nnn { grid / lines } { pitch }    { document-grid / lines / pitch }
\__hit_layout_leaf:nnn { grid / lines } { per-page } { document-grid / lines / lines-per-page }
%    \end{macrocode}
%
% \option{grid}/\option{chars}/\option{latin} 决定网格字距加不加到西文字母上。
% Word 那边 \texttt{charSpace} 是加在\emph{每一个}字符后面的，汉字、字母、数字、
% 空格一视同仁——实测每个西文字母后面也多 0.4543\,bp，一行 82 字累计 37\,bp，
% 差不多三个汉字宽。默认关，因为它会动到参考文献、URL 与代码环境里的每一个
% 西文串。
%
% 原先叫 \option{char-pitch-en}，那个名字像是有个 \option{char-pitch-zh} 与它配对，
% 其实没有：汉字那一侧不是全文开关，是 \option{styles} 里各样式
% \option{font} 下的 \option{snap-to-grid}，按元素发 \cs{CJKglue}。两边不对称是
% 实现逼出来的——\cs{CJKglue} 是个宏，能按组换；西文那边要发
% \cs{defaultfontfeatures} 的 \texttt{LetterSpace}，得重新声明整个字族，
% 只能全文来一次。所以它留在 \option{layout}。
%    \begin{macrocode}
\bool_new:N \g__hit_layout_pitch_en_bool
\__hit_boolean_key:nnn { layout / grid / chars } { latin } { layout_pitch_en }
\__hit_boolean_key:nnn { layout } { grid / chars / latin } { layout_pitch_en }
%    \end{macrocode}
%
% \option{grid}/\option{enabled} 管这一节的网格启不启用。Word 那边是
% \texttt{docGrid} 有没有 \texttt{w:type}：写了 \texttt{linePitch} 与
% \texttt{charSpace} 却没写 \texttt{w:type}，那两个数就只是记着，一个都不生效。
% 深圳本科报告的正文是这一档，本部本科那两份也是。
%
% 它管两件事。一是贴不贴网格，与 \option{styles} 里逐目标写
% \option{snap-to-grid} 是一回事，这里是全节一刀切。二是\emph{「一行」这个单位
% 折成多长}：网格启用时是网格行距，关着时是这一节正文的单倍行高，
% 也就是 \option{styles}/\option{Normal} 那一档的字号。
%
% \emph{第二件事不能省。} 章节标题的段前段后写的是「1 行」「0.8 行」「0.5 行」，
% 网格关着时 Word 折成正文的一个字号，深圳本科那两份原件段落上缓存的正是
% 240／192／120 缇即 12／9.6／6\,bp；按网格行距折会得到 19.15／15.32／9.58，
% 每个章标题差七个 bp，一页下来越积越多。
%    \begin{macrocode}
\bool_new:N \g__hit_layout_grid_bool
\bool_gset_true:N \g__hit_layout_grid_bool
\__hit_boolean_key:nnn { layout / grid } { enabled } { layout_grid }
\__hit_boolean_key:nnn { layout } { grid / enabled } { layout_grid }
%    \end{macrocode}
%
%    \begin{macrocode}
\bool_new:N \g__hit_layout_frozen_bool
\AddToHook { begindocument } { \bool_gset_true:N \g__hit_layout_frozen_bool }
\cs_new_protected:Npn \hit@layout@refresh:
  {
    \bool_if:NF \g__hit_layout_frozen_bool
      {
        \__hit_layout_msword_active:T
          {
            \tl_if_empty:NF \g__hit_layout_msword_tl
              { \exp_args:NV \docxlike_page_setup:n \g__hit_layout_msword_tl }
          }
      }
  }
\cs_new_protected:Npn \hit@layout@set:n #1
  {
    \bool_set_false:N \l__hit_layout_dirty_bool
    \hit@keys@set:nn { hit / layout } {#1}
    \bool_if:NT \l__hit_layout_dirty_bool
      { \exp_args:NV \docxlike_page_setup:n \g__hit_layout_msword_tl }
  }
%    \end{macrocode}
%
% \pkg{etoolbox} 在这里装。下面 \cs{AfterPreamble} 与 \cs{ifdefempty} 要它，
% \file{src/engine/hit-caption.dtx} 的三处 \cs{patchcmd} 也要。两个消费者都排在
% \cs{LoadClass} 之前，\file{src/flow/hit-*-deps.dtx} 太晚，所以在这儿。
% \pkg{etoolbox} 不牵扯字体与基类，提前装没有副作用。
%    \begin{macrocode}
\RequirePackage{etoolbox}
%</shared-keylist>
%    \end{macrocode}
%
% \section{终稿的键清单（keylist）}
%
% 学位论文用到的元信息字段：字段名、旧名映射、\cs{hitsetup} 各族的路由，
% 以及摘要环境。都是清单，不含机制。
%
% ^^A ======================================================================
% ^^A Module thesis-keylist: 元信息字段（\hitsetup 接受的 key）（hit-thesis）
% ^^A ======================================================================
%    \begin{macrocode}
%<*shared-keylist>

\hit@define@term{secret-level} %密级
\hit@define@term{classified-index-national}  %国内图书分类号
\hit@define@term{classified-index-international}  %国际图书分类号

\hit@define@term{discipline-zh} %中文学科
\hit@define@term{associate-supervisor-zh} %中文副导师
\hit@define@term{co-supervisor-zh}%中文联合导师
%    \end{macrocode}
%
% \changes{v3.1b}{2022/06/06}{增加深圳封面最下方的时间信息}
% \changes{v3.1d}{2025/03/03}{根据 \issue[226]，将 \texttt{szshortcdate} 改为 \texttt{firstpagecdate}，只要赋值就会显示}
% 根据 \issue[95] 增加一个变量 \texttt{szshortcdate} 来存储答辩时间。
% 根据 \issue[226]，将 \texttt{szshortcdate} 改为 \texttt{firstpagecdate}，只要赋值就会显示。
%    \begin{macrocode}
\hit@define@date@term{cover-date-zh}{zh}
\hit@define@term{report-number}
\hit@define@term{station-zh}
%    \end{macrocode}
%
% 添加博后封皮选项
% \changes{v3.0.12}{2020/11/02}{添加博后封皮选项}
%    \begin{macrocode}
\hit@define@date@term{work-finish-date-zh}{zh}
\hit@define@date@term{report-submit-date-zh}{zh}
\hit@define@date@term{research-start-date-zh}{zh}
\hit@define@date@term{research-end-date-zh}{zh}
%    \end{macrocode}
%
% 添加博后封皮的title选项
% \changes{v3.0.0}{2020/05/21}{添加博后封皮的title选项}
%    \begin{macrocode}




\hit@define@title{en}
\hit@define@term{discipline-en} %英文学科
\hit@define@term{author-en} %英文作者
\hit@define@term{supervisor-en} %英文导师
\hit@define@term{associate-supervisor-en} %英文副导师
\hit@define@term{co-supervisor-en} %英文联合导师
\hit@define@term{affiliation-en}
\hit@define@term{major-en}
\hit@define@date@term{date-en}{en}
%    \end{macrocode}
%
% 摘要是「先收着，后面再排」：用户在正文位置写 \env{abstract-zh}，内容存进
% \cs{hit@abstract@zh}，真正排版在摘要页那边。所以要把整个环境体当参数抓下来。
% \cs{NewDocumentEnvironment} 的 \texttt{b} 参数类型就是环境体，必须排在最后，
% 前面的 \texttt{+} 表示允许 \cs{par}。
%
% 环境体是当 token 抓的，里面不能放 \cs{verb} 或 \env{lstlisting}。
% 原先的 \cs{Collect@Body} 是一样的限制。
%    \begin{macrocode}
\cs_new_protected:Npn \hit@collect@abstract@zh #1 { \cs_gset:Npn \hit@abstract@zh {#1} }
\cs_new_protected:Npn \hit@collect@abstract@en #1 { \cs_gset:Npn \hit@abstract@en {#1} }
\NewDocumentEnvironment { abstract-zh } { +b } { \hit@collect@abstract@zh {#1} } { }
\NewDocumentEnvironment { abstract-en } { +b } { \hit@collect@abstract@en {#1} } { }
\NewDocumentEnvironment { cabstract } { +b }
  {
    \hit@warn@renamed@environment { cabstract } { abstract-zh }
    \hit@collect@abstract@zh {#1}
  }
  { }
\NewDocumentEnvironment { eabstract } { +b }
  {
    \hit@warn@renamed@environment { eabstract } { abstract-en }
    \hit@collect@abstract@en {#1}
  }
  { }
%    \end{macrocode}
%
% 关键词的拆分机制见 \file{src/engine/hit-key-engine.dtx}。
%    \begin{macrocode}
\hit@parse@keywords{keywords-zh}{abstract@keywords@separator@zh}
\hit@parse@keywords{keywords-en}{abstract@keywords@separator@en}
%    \end{macrocode}
%
% \cs{hitsetup} 的入口换成 l3keys。根族只做路由，把 \option{info}、\option{style}
% 等分派到对应子族；子族一律把键原样转回 \texttt{hit/class}
% 族，所以两种写法结果相同：
% \begin{latex}
% \hitsetup{title-zh={...}}              % 扁平写法，一直支持
% \hitsetup{info={title-zh={...}}}       % 新的分族写法
% \end{latex}
% 根族的 \option{unknown} 兜底负责扁平写法，子族的兜底负责分族写法。等某个键真正
% 迁到 l3keys 之后，它就不再走兜底，两条路自然收敛。兜底本身、三态取值与布尔键的
% 声明机制都在 \file{src/engine/hit-key-engine.dtx}。
%
% 元信息字段名单。这批由 \cs{hit@define@term} 与 \cs{hit@parse@keywords} 批量生成，
% 每个字段同时有一个同名命令（\cs{title-zh} 等），键的处理器直接调用那个命令，
% 所以三种写法结果一致：\cs{hitsetup} 的扁平写法、分族写法、以及命令写法。
%    \begin{macrocode}
\hit@declare@info@field
  {
    secret-level , classified-index-national , classified-index-international ,
    discipline-zh , associate-supervisor-zh , co-supervisor-zh , cover-date-zh ,
    report-number , station-zh , work-finish-date-zh , report-submit-date-zh ,
    research-start-date-zh , research-end-date-zh , title-en ,
    discipline-en , author-en , supervisor-en , associate-supervisor-en ,
    co-supervisor-en , affiliation-en , major-en , date-en , keywords-zh , keywords-en
  }
%    \end{macrocode}
%
% 可本地化字符串的名单。默认值在类文件末尾给出，\cs{hitsetup} 在导言区执行，
% 所以用户设的值覆盖默认值。键名怎么由宏名得来见 \file{src/engine/hit-key-engine.dtx}。
%    \begin{macrocode}
\hit@declare@constant
  {
    declarations-statement-bachelor , declarations-authorization-bachelor ,
    declarations-heading-line-one-bachelor , declarations-heading-line-two-bachelor ,
    declarations-heading-short-bachelor , acknowledgements-heading-zh ,
    acknowledgements-heading-en , declarations-heading-zh , declarations-heading-en ,
    declarations-authorization-text , declarations-statement-text-bachelor , bibliography-heading-en ,
    base-brace-left-zh , base-brace-right-zh , abstract-keywords-separator-zh , abstract-keywords-zh ,
    conclusion-heading-zh , abstract-keywords-en , table-of-contents-page-en ,
    table-of-contents-table-en , table-of-contents-figure-en , conclusion-heading-en ,
    table-of-contents-undefined-en ,
    resume-corresponding-address-zh , resume-corresponding-address-en ,
    declarations-statement , declarations-heading-bachelor-zh , declarations-statement-text ,
    degree-type-zh , degree-type-en ,
    degree-type-academic-zh , degree-type-academic-en ,
    degree-type-professional-zh , degree-type-professional-en ,
    list-of-symbols-heading-zh , list-of-symbols-heading-en ,
    list-of-abbreviations-heading-zh , list-of-abbreviations-heading-en ,
    list-of-symbols-and-abbreviations-heading-zh , list-of-symbols-and-abbreviations-heading-en ,
    publications-doctoral-heading-zh ,
    publications-doctoral-heading-en , titlepage-affiliation-en , titlepage-associate-supervisor-en , titlepage-author-en , base-brace-left-en ,
    base-brace-right-en , titlepage-co-supervisor-en , titlepage-date-en , base-degree-of-discipline-en , titlepage-degree-applied-en , abstract-keywords-separator-en ,
    base-university-en , titlepage-institution-en , titlepage-major-en , titlepage-supervisor-en , base-degree-level-en , titlepage-degree-attribute , base-heading-spread ,
    index-heading-en , index-heading-zh , titlepage-classified-index-international-en ,
    titlepage-classified-index-national-en , cover-udc-postdoc ,
    titlepage-author-postdoc-zh , cover-classified-index-postdoc , base-report-postdoc ,
    titlepage-research-end-date-postdoc , cover-work-finish-date-postdoc ,
    titlepage-major-postdoc-zh , cover-number-postdoc ,
    titlepage-station-postdoc , base-secrecy-postdoc , titlepage-separator-postdoc ,
    titlepage-research-start-date-postdoc , cover-report-submit-date-postdoc ,
    titlepage-co-supervisor-postdoc-zh , publications-heading-zh , publications-heading-en ,
    resolution-heading-zh , resolution-heading-en , resume-heading-zh , resume-heading-en ,
    resolution-reviewer , resolution-name , resolution-title ,
    resolution-affiliation , resolution-discipline , resolution-role ,
    resolution-committee , resolution-chair , resolution-member ,
    resolution-secretary , resolution-text ,
    declarations-authorization-text-bachelor , cover-title-label-bachelor , base-title-separator-en
  }
%    \end{macrocode}
%
% 两个族的键。排版开关都是普通布尔，用 \cs{\_\_hit\_style\_pair:nn} 成对声明，
% 新名设的是旧标志位本身，所以两套名字天然一致。
% \option{heading-latin-sans} 跟共享层的 \option{toc}/\option{latin-sans} 是一对，
% 一个管章节标题里的西文一个管目录里的。参考文献的五个键此前分散在
% \file{book} 与 \file{art} 两个模块、命名也不统一（\texttt{citetwo} 看不出是
% 控制什么的），新旧两套名字设的是同一批内部变量，所以可以混用。
%    \begin{macrocode}
%    \end{macrocode}
%
% \option{layout}/\option{openright} 决定强制另起一页时要不要落在奇数页。
% 写在全篇一级是默认值，写在部件名下只管那个部件：
%
% \begin{quote}
% \cs{hitsetup}\{ layout = \{ openright = false ,
%     cover = \{ openright = true \} \} \}
% \end{quote}
%
% 部件那一层收 \texttt{auto}，含义是「我不表态，用全篇的值」。交图书馆的电子版
% 要求没有任何空白页，写 \option{layout}\texttt{=\{openright=false\}} 即可，
% 不必另设开关。
%
% \emph{为什么在 \option{layout} 里。} Word 的 \texttt{w:sectPr} 同时装着页边距
% 与分节符类型，页面几何与页流本来就是一层的事；我们的 \option{layout} 建模的
% 正是它。
%    \begin{macrocode}
\clist_map_inline:nn
  {
  }
  { \__hit_style_pair:nn #1 }
%    \end{macrocode}
%
% \option{style}/\option{em-dash} 决定破折号（\texttt{U+2014}）从哪一副字取字形，
% 取 \texttt{cjk} 或 \texttt{latin}，默认 \texttt{cjk}。
%
% \emph{两副字的杠长得不一样。} 中易宋体的占 0.031\,--\,0.973\,em，两个「\,——\,」
% 中间断开一截；\pkg{Times} 的占 $-0.009$\,--\,1.009\,em，两头各出头一点、连成一条，
% 只是位置低些。Word 里这个字看 \texttt{run} 上的 \texttt{w:hint}：写着
% \texttt{eastAsia} 就排中文字，没写就排西文字，全凭作者的输入法留下什么。
%
% 取值按变体给，见 \file{hithesis.cfg}。判据是逐份数原件里的 \texttt{hint}：
% 深圳本科中期那份 \file{docx} 里 11 个破折号 \texttt{run} 一个
% \texttt{eastAsia} 都没有，所以深圳的报告一律西文；学位论文那两份范例
% \texttt{eastAsia} 占多数，仍是中文。
%
% 引号、省略号、间隔号不归它管，照旧从中文字体取。
%    \begin{macrocode}
\bool_new:N \g__hit_em_dash_latin_bool
\keys_define:nn { hit / style }
  {
    em-dash .choice: ,
    em-dash / cjk .code:n =
      { \bool_gset_false:N \g__hit_em_dash_latin_bool } ,
    em-dash / latin .code:n =
      { \bool_gset_true:N \g__hit_em_dash_latin_bool } ,
    em-dash .value_required:n = true ,
  }
%    \end{macrocode}
%
% \option{style}/\option{fake-slant-songti} 管宋体的斜体，布尔，默认
% \texttt{false}。汉字没有斜体字面，宋体的斜体默认换楷体（写作指南与
% \pkg{ctex} 都这么做）；开了就把宋体的常规字面斜切，全文档凡是宋体的地方
% （正文字族与 \cs{songti}）一律如此，楷体只由 \cs{kaishu} 选用。别的字族
% （黑体、楷体、仿宋）没有这个选择，斜体一律斜切，不归它管。
% 倾斜因子与落实的办法见 \file{src/engine/hit-fonts.dtx} 的 \cs{hit@font@italic@pass:}。
%    \begin{macrocode}
\bool_new:N \g__hit_font_fake_slant_songti_bool
\keys_define:nn { hit / style }
  {
    fake-slant-songti .bool_gset:N = \g__hit_font_fake_slant_songti_bool ,
    fake-slant-songti .default:n = { true } ,
  }
%    \end{macrocode}
%
% \option{style}/\option{verb-space} 管行内代码（\cs{verb}）里的空格印成什么，取
% \texttt{visible}（默认，U+02FD 那个开口方框）或 \texttt{cell}（不可见的不断行空格
% U+00A0）。两档都是一格宽、刚性、可在其后断行。排法与理由见
% \file{src/docxlike/docxlike-format.dtx} 的「行内代码」一节。
%    \begin{macrocode}
\bool_new:N \g__hit_verb_space_visible_bool
\bool_gset_true:N \g__hit_verb_space_visible_bool
\keys_define:nn { hit / style }
  {
    verb-space .choice: ,
    verb-space / visible .code:n =
      { \bool_gset_true:N \g__hit_verb_space_visible_bool \hit@docxlike@sync: } ,
    verb-space / cell .code:n =
      { \bool_gset_false:N \g__hit_verb_space_visible_bool \hit@docxlike@sync: } ,
    verb-space .value_required:n = true ,
  }
%    \end{macrocode}
%
% \option{style}/\option{float-sep} 是浮动体（图、表）与前后正文之间的空当，
% 收 \texttt{auto} 或一个长度，默认 \texttt{auto}。\texttt{auto} 照各档原来的
% 规矩：学位论文按指南 2.12「插表的上下与文中文字间需空一行编排」给一个网格行
% （\file{src/flow/hit-final-float.dtx}），报告给 \pkg{ccaption} 那一档的 12\,bp
% （\file{src/engine/hit-caption.dtx}）。写长度就三个位置（正文中间、页顶页底、
% 两个浮动体之间）一律这么多，不给伸缩。深圳本科的两份报告在 \file{hithesis.cfg}
% 里写 \texttt{0bp}：原件里图段表段一个 \texttt{w:before}／\texttt{w:after} 都
% 没有，空一行是作者自己敲的空段。
%    \begin{macrocode}
\tl_new:N \g__hit_float_sep_tl
\keys_define:nn { hit / style }
  {
    float-sep .code:n =
      {
        \str_if_eq:nnTF {#1} { auto }
          { \tl_gclear:N \g__hit_float_sep_tl }
          { \tl_gset:Nx \g__hit_float_sep_tl { \dim_eval:n {#1} } }
      } ,
    float-sep .value_required:n = true ,
  }
%    \end{macrocode}
%
% \texttt{math} 族。\option{numbering-parenthesis-style} 取 \texttt{full} 是全角
% 括号、\texttt{half} 是半角，默认半角——本科范例印出来就是半角，v3.1e 也是。
% 英文论文不受它影响，一律半角，那个判断在 \cs{hit@equation@paren@left} 里。
%    \begin{macrocode}
\keys_define:nn { hit / math }
  {
    numbering-parenthesis-style .choice: ,
    numbering-parenthesis-style / full .code:n =
      { \bool_gset_true:N \g__hit_equation_paren_full_bool } ,
    numbering-parenthesis-style / half .code:n =
      { \bool_gset_false:N \g__hit_equation_paren_full_bool } ,
    numbering-parenthesis-style / auto .code:n = { } ,
    numbering-parenthesis-style .value_required:n = true ,
  }
%    \end{macrocode}
%
% \option{compress-min} 是「相邻连续序号几个起压成区间」：\texttt{2} 排
% \texttt{[1-2]}（默认），\texttt{3} 排 \texttt{[1,2]}，三个起总归压。两条线
% 都认，落点都是官方机制：传统一线转发给 \pkg{gbt7714.sty} 的
% \texttt{mincompress}（v3.2a 起，原先自家打的 \pkg{natbib} 补丁退役），
% biblatex 一线落到 \pkg{gb7714-2015} 的 \texttt{gbrefcompress} 计数器。
% 两边都不锁时机，随设随生效；旧类选项 \option{citetwo} 的那对标志位并进
% 同一个取值函数，两条线都认它。
%
% \option{uppercase-family}（\texttt{true} 著者姓名全大写即国标排法、
% \texttt{false} 保持 \file{.bib} 原样）、\option{max-bib-names}（文献表里
% 最多列几位作者，\texttt{0} 不限）、\option{max-cite-names}（正文引用侧的
% 同款，数字引用不出姓名，只在 biblatex 一线对 \cs{textcite} 这类起作用）、
% \option{numbering-align}（标号 \texttt{left}/\texttt{right}）这几个在
% biblatex 一线是装载时换算成它的选项的，装载之后再设不生效，设值器前面搭
% 一句守卫（见 \file{hit-bibstyle.dtx}）。\option{split-entry} 排版时才读，
% 不用守。
%    \begin{macrocode}
\int_new:N \g__hit_cite_compress_int
\cs_new:Npn \__hit_bib_compress_effective:
  {
    \int_compare:nNnTF \g__hit_cite_compress_int > 0
      { \int_use:N \g__hit_cite_compress_int }
      {
        \bool_if:NTF \g__hit_citetwo_comma_bool { 3 }
          { \bool_if:NTF \g__hit_citetwo_endash_bool { 2 } { 0 } }
      }
  }
\cs_new_protected:Npn \hit@bib@cite@compress@apply:
  {
    \int_compare:nNnF { \__hit_bib_compress_effective: } = 0
      {
        \@ifpackageloaded { biblatex }
          {
            \cs_if_exist:cT { c@gbrefcompress }
              { \setcounter { gbrefcompress } { \__hit_bib_compress_effective: } }
          }
          {
            \@ifpackageloaded { gbt7714 }
              {
                \__hit_bib_if_mincompress:T
                  {
                    \use:e
                      {
                        \exp_not:N \keys_set:nn { gbt7714 }
                          { mincompress = \__hit_bib_compress_effective: }
                      }
                  }
              } { }
          }
      }
  }
%    \end{macrocode}
%
% \texttt{mincompress} 是 \pkg{gbt7714}~v3 才有的键。\TeX~Live 2025 装的是
% v2.1.5（2022/10/03），\TeX~Live 2026 才换成 v3.0.1（2026/07/20）。直接
% \cs{keys\_set:nn} 过去，v2 上会报 \texttt{The key 'gbt7714/mincompress' is
% unknown}，配上 \texttt{-halt-on-error} 就是编不出来。
%
% 判据用 \cs{keys\_if\_exist:nnTF} 查键本身，不查包的版本号或日期：版本号将来
% 怎么走不好说，键在不在才是我们真正依赖的东西。这个判断在 \TeX~Live 2025 与
% 2026 的 \pkg{l3kernel} 里都有；万一更老的树上没有，外面那层
% \cs{cs\_if\_exist:NTF} 会让它退回「不转发」，仍然编得出来，只是压缩阈值用
% \pkg{gbt7714} 自己的默认值。
%    \begin{macrocode}
\prg_new_protected_conditional:Npnn \__hit_bib_if_mincompress: { T }
  {
    \cs_if_exist:NTF \keys_if_exist:nnTF
      {
        \keys_if_exist:nnTF { gbt7714 } { mincompress }
          { \prg_return_true: } { \prg_return_false: }
      }
      { \prg_return_false: }
  }
%    \end{macrocode}
%
%    \begin{macrocode}
\keys_define:nn { hit / bib }
  {
    compress-min .choices:nn = { auto , 2 , 3 }
      {
        \str_case:Vn \l_keys_choice_tl
          {
            { auto } { \int_gset:Nn \g__hit_cite_compress_int { 0 } }
            { 2 }
              {
                \int_gset:Nn \g__hit_cite_compress_int { 2 }
                \bool_gset_false:N \g__hit_citetwo_comma_bool
                \bool_gset_true:N  \g__hit_citetwo_endash_bool
              }
            { 3 }
              {
                \int_gset:Nn \g__hit_cite_compress_int { 3 }
                \bool_gset_true:N  \g__hit_citetwo_comma_bool
                \bool_gset_false:N \g__hit_citetwo_endash_bool
              }
          }
        \hit@bib@cite@compress@apply:
      } ,
    compress-min .value_required:n = true ,
    numbering-align .code:n =
      {
        \str_if_eq:nnF {#1} { auto }
          {
            \hit@bib@key@guard:n { numbering-align }
            \str_gset:Nn \g_bib_labelalign_str {#1}
          }
      } ,
    uppercase-family .choices:nn = { true , false , auto }
      {
        \str_if_eq:VnF \l_keys_choice_tl { auto }
          {
            \hit@bib@key@guard:n { uppercase-family }
            \str_if_eq:VnTF \l_keys_choice_tl { true }
              { \bool_gset_true:N  \g_bib_authorupper_bool }
              { \bool_gset_false:N \g_bib_authorupper_bool }
          }
      } ,
    uppercase-family .default:n = true ,
    max-bib-names .code:n =
      {
        \str_if_eq:nnF {#1} { auto }
          {
            \hit@bib@key@guard:n { max-bib-names }
            \int_gset:Nn \g_bib_maxauthor_int {#1}
            \int_if_zero:nT { \g_bib_maxauthor_int }
              { \int_gdecr:N \g_bib_maxauthor_int }
          }
      } ,
    max-cite-names .code:n =
      {
        \str_if_eq:nnF {#1} { auto }
          {
            \hit@bib@key@guard:n { max-cite-names }
            \int_gset:Nn \g__hit_bib_maxcite_int {#1}
            \int_compare:nNnT \g__hit_bib_maxcite_int = 0
              { \int_gdecr:N \g__hit_bib_maxcite_int }
          }
      } ,
    split-entry .choice: ,
    split-entry / true  .code:n = { \bool_gset_true:N \g__hit_splitbibitem_bool } ,
    split-entry / false .code:n = { \bool_gset_false:N \g__hit_splitbibitem_bool } ,
    split-entry / auto  .code:n = { } ,
    split-entry .default:n = true ,
  }
%    \end{macrocode}
%
% 兼容期的处理。旧名一律照常生效，只在日志里给一条提示。三类提示分开：
% 改过名的选项逐个报，元信息字段的扁平写法合并成一条报，
% \cs{documentclass} 里传排版类选项单独报。
%
% 另有一张废弃表，记的是没有新名可改的选项。这类键照旧认，只是什么也不做，
% 提示里说清为什么。\option{tocblank} 插的那个空行规范里没有，所以整条去掉，
% 行为等同于原先的 \option{tocblank}\texttt{=false}。
% \option{glue} 和 \option{raggedbottom} 不能只报一个新名字。它们原来同时
% 承担两件事：调旧版面那套手写的伸缩量，以及决定页底排不排齐。新版面把这
% 两件事分开了，所以得问清楚用户到底要哪一种，两条路都写给他看。
%    \begin{macrocode}
\prop_gset_from_keyval:Nn \g__hit_split_option_prop
  {
    glue         = { split } ,
    raggedbottom = { split } ,
  }
%    \end{macrocode}
%
%    \begin{macrocode}
\prop_gset_from_keyval:Nn \g__hit_removed_option_prop
  {
    tocblank = { 目录里第一章之前不再插空行，规范里没有这一条 } ,
  }
\prop_gset_from_keyval:Nn \g__hit_renamed_option_prop
  {
    arialtoc         = { toc = { latin-sans = ... } } ,
    tocfour          = { toc = { depth = 4 } } ,
    engtoc           = { toc = { english = ... } } ,
    arialtitle       = { styles = { heading = { font = { latin-text-font = Arial } } } } ,
    chapterbold      = { styles = { Heading 1 = { font = { font-style = bold } } } } ,
    chapterhang      = { styles = { Heading 1 = { paragraph = { indentation = { special = hanging } } } } } ,
    capcenterlast    = { styles = { caption = { paragraph = { alignment = last-line-centered } } } } ,
    subcapcenterlast = { styles = { caption-sub = { paragraph = { alignment = last-line-centered } } } } ,
    absupper         = { styles = { abstract-heading = { font = { all-caps = ... } } } } ,
    etitleauto       = { titlepage = { title-en-xiaoer = ... } } ,
    openright        = { layout = { openright = ... } } ,
    library          = { layout = { openright = false } } ,
    zijv             = { layout = v2-layout~专用 } ,
    bsheadrule       = { layout = { header = { rule = ... } } } ,
    bsmainpagenumberline  = { layout = { mainmatter = { page-number-line = ... } } } ,
    bsfrontpagenumberline = { layout = { frontmatter = { page-number-line = ... } } } ,
    citetwo          = { bib = { compress-min = 2~|~3 } } ,
    splitbibitem     = { bib = { split-entry = ... } } ,
    biblabelalign    = { bib = { numbering-align = ... } } ,
    bibauthorupper   = { bib = { uppercase-family = true } } ,
    bibmaxauthor     = { bib = { max-bib-names = ... } } ,
    raggedbottom     = { layout = { bottom = ragged~|~flush } } ,
  }
%    \end{macrocode}
%
% 旧名到新名的对照。旧名在 v3.1e 就是公开接口，写在用户的 \file{cover.tex} 里，
% 所以保留到 v4。
%    \begin{macrocode}
\hit@define@legacy@constant{cstudenttype}{degree-type-zh}
\hit@define@legacy@constant{estudenttype}{degree-type-en}
\hit@define@legacy@terms
  {
    statesecrets=secret-level ,
    natclassifiedindex=classified-index-national ,
    intclassifiedindex=classified-index-international ,
    ctitle=title-zh ,
    cxueke=discipline-zh ,
    cauthor=author-zh ,
    csupervisor=supervisor-zh ,
    cassosupervisor=associate-supervisor-zh ,
    ccosupervisor=co-supervisor-zh ,
    caffil=affiliation-zh ,
    csubject=major-zh ,
    firstpagecdate=cover-date-zh ,
    cdate=date-zh ,
    cnumber=report-number ,
    cpositionname=station-zh ,
    cfinishdate=work-finish-date-zh ,
    csubmitdate=report-submit-date-zh ,
    cstartdate=research-start-date-zh ,
    cenddate=research-end-date-zh ,
    cstudentid=student-id ,
    etitle=title-en ,
    exueke=discipline-en ,
    eauthor=author-en ,
    esupervisor=supervisor-en ,
    eassosupervisor=associate-supervisor-en ,
    ecosupervisor=co-supervisor-en ,
    eaffil=affiliation-en ,
    esubject=major-en ,
    edate=date-en ,
    ckeywords=keywords-zh ,
    ekeywords=keywords-en
  }
\hit@define@legacy@title{ctitlecover}{zh}{first}
\hit@define@legacy@title{ctitleone}{zh}{first}
\hit@define@legacy@title{ctitletwo}{zh}{second}
\hit@define@legacy@title{csubtitle}{zh}{subtitle}
\hit@define@legacy@title{esubtitle}{en}{subtitle}
\prop_gset_from_keyval:Nn \g__hit_renamed_load_option_prop
  {
    font        = { font = { latin-text-font = ... } } ,
    cjk-font    = { font = { asian-text-font = ... } } ,
    type        = { degree-level } ,
    newgeometry = { v2-layout~=~{~bachelor~,~graduate~=~one|two~} } ,
  }
%    \end{macrocode}
%
% 单双面由学位与校区决定，所以扫描类选项时额外查一下 \option{oneside} 与
% \option{twoside}。另外两个类没有这条。
%    \begin{macrocode}
\cs_set_protected:Npn \__hit_class_option_extra:
  {
    \clist_if_in:nVT { oneside , twoside } \l__hit_warn_message_tl
      {
        \exp_args:NV \__hit_warn_once:nnn \l__hit_warn_message_tl
          { manual-sides } { }
      }
  }
%    \end{macrocode}
%
% \texttt{defense} 族收答辩决议那一页要的数据。一个人一条记录，字段具名：
% \begin{latex}
% \hitsetup{ defense = {
%   reviewer = {name=张三, title=教授（博导）, affiliation=哈尔滨工业大学,
%               discipline=控制科学与工程},
%   reviewer = {name=李四, …},
%   chair     = {…},
%   member    = {…},
%   secretary = {…},
%   resolution = {答辩委员会经过讨论，一致同意通过答辩，建议授予工学博士学位。},
% } }
% \end{latex}
%
% \textbf{同一个键写多遍是累加，不是覆盖。} \option{reviewer} 与 \option{member}
% 靠这个收任意条。\pkg{l3keys} 按书写顺序逐条执行，键的代码是「往队列右边追加」，
% 写几遍就是几个人，顺序也就是写的顺序。
%
% 四个名单各一个队列，与纸质表的分块一一对应。每一项存的是那个人的整串键值，
% 排版时再解开，所以少写一个字段只是那一格空着，不会像并列列表那样把后面全顶错位。
%    \begin{macrocode}
\seq_new:N \g__hit_defense_reviewer_seq
\seq_new:N \g__hit_defense_chair_seq
\seq_new:N \g__hit_defense_member_seq
\seq_new:N \g__hit_defense_secretary_seq
\tl_new:N  \g__hit_defense_resolution_tl
%    \end{macrocode}
%
% \subsection{缩略语词条}
%
% \option{abbr} 族收缩略语词条，形状与 \option{defense} 一样——数据进族，
% 部件只管排不排。键名就是短形，值是那一条的字段：
% \begin{latex}
% \hitsetup{ abbr = {
%   SCNA = {zh=体细胞拷贝数变异, en=Somatic copy number alteration},
%   SVM  = {Support Vector Machine},
% } }
% \end{latex}
% 值不含等号时按裸值处理，含汉字归 \option{zh}，不含归 \option{en}。
% \option{label} 改标签，不写就取短形。
%
% 另有两个开关：\option{indexing} 决定 \cs{abbr} 要不要顺带登记索引，
% \option{reset-after-abstract} 决定摘要排完要不要把「用过」状态清空，
% 让正文重新算首次出现。两个默认都开。
%    \begin{macrocode}
\str_new:N \g__hit_abbr_punct_str
\str_gset:Nn \g__hit_abbr_punct_str { auto }
\bool_new:N \g__hit_abbr_indexing_bool
\bool_new:N \g__hit_abbr_reset_bool
\bool_gset_true:N \g__hit_abbr_indexing_bool
\bool_gset_true:N \g__hit_abbr_reset_bool
\keys_define:nn { hit / abbr }
  {
    indexing .bool_gset:N = \g__hit_abbr_indexing_bool ,
    indexing .default:n = true ,
%    \end{macrocode}
%
% \option{punct} 管首次出现那对括号与里头的逗号取全角还是半角。
% \texttt{auto}（默认）按语种走：中文版全角、英文版半角。中文数学类论文里
% 行文用半角标点的也不少，那时候写 \texttt{half} 一次改掉。
%    \begin{macrocode}
    punct .choices:nn = { auto , full , half }
      { \str_gset:NV \g__hit_abbr_punct_str \l_keys_choice_tl } ,
    punct .value_required:n = true ,
    reset-after-abstract .bool_gset:N = \g__hit_abbr_reset_bool ,
    reset-after-abstract .default:n = true ,
    unknown .code:n =
      { \exp_args:NV \hit@abbr@declare:nn \l_keys_key_str {#1} } ,
  }
%    \end{macrocode}
%
%    \begin{macrocode}
\keys_define:nn { hit / defense }
  {
    reviewer   .code:n = { \seq_gput_right:Nn \g__hit_defense_reviewer_seq  {#1} } ,
    chair      .code:n = { \seq_gput_right:Nn \g__hit_defense_chair_seq     {#1} } ,
    member     .code:n = { \seq_gput_right:Nn \g__hit_defense_member_seq    {#1} } ,
    secretary  .code:n = { \seq_gput_right:Nn \g__hit_defense_secretary_seq {#1} } ,
    resolution .code:n = { \tl_gset:Nn \g__hit_defense_resolution_tl {#1} } ,
    unknown .code:n =
      { \msg_warning:nnV { hithesis } { unknown-key } \l_keys_key_str } ,
  }
%    \end{macrocode}
%
% 一条记录里的四个字段。排版时把一项喂进这个族，四个变量就有了值；
% 每次解开之前先清空，免得上一个人的值漏到下一个人身上。
%    \begin{macrocode}
\tl_new:N \l__hit_defense_name_tl
\tl_new:N \l__hit_defense_title_tl
\tl_new:N \l__hit_defense_affiliation_tl
\tl_new:N \l__hit_defense_discipline_tl
\keys_define:nn { hit / defense / person }
  {
    name        .tl_set:N = \l__hit_defense_name_tl ,
    title       .tl_set:N = \l__hit_defense_title_tl ,
    affiliation .tl_set:N = \l__hit_defense_affiliation_tl ,
    discipline  .tl_set:N = \l__hit_defense_discipline_tl ,
    unknown .code:n =
      { \msg_warning:nnV { hithesis } { unknown-key } \l_keys_key_str } ,
  }
\cs_new_protected:Npn \hit@defense@unpack:n #1
  {
    \tl_clear:N \l__hit_defense_name_tl
    \tl_clear:N \l__hit_defense_title_tl
    \tl_clear:N \l__hit_defense_affiliation_tl
    \tl_clear:N \l__hit_defense_discipline_tl
    \hit@keys@set:nn { hit / defense / person } {#1}
  }
%    \end{macrocode}
%    \begin{macrocode}
\__hit_check_class_options:
%    \end{macrocode}
%
% 封面一族。
%
% \option{student-type} 决定封面上排不排学位类别那一行。\texttt{auto} 按校区与
% 学位判：深圳、硕士、非全日制这三样占一样就排，本部全日制博士不排。给
% \texttt{true} 或 \texttt{false} 就照给的来，也收
% \texttt{\{bachelor=…,master=…,doctor=…\}} 这种按学位分开写的形式。
%
% 这条判据是 \issue[97] 与 \issue[149] 两次来回改出来的，当时议的是深圳硕士要不要
% 加，顺手把本部全日制博士排除了；而研究生书写范例的封面本身就是本部全日制博士，
% 上面是有这一行的。判据与范例对不上，所以留一个显式开关，不动默认。
%
% \option{title-en-xiaoer} 把英文题目强制排小二号。不开时字号由封面
% 自己按塞不塞得下定：先按二号塞，空行让步全用尽还装不下才降小二——规范给的是
% 二号，「题目太长时可用小二号」，降号是允许的一档但不为了排得好看去缩号。
% 想还没到极限就用小二，把这一条开起来。
%    \begin{macrocode}
\tl_new:N \g_hit_cover_title_en_xiaoer_tl
\tl_gset:Nn \g_hit_cover_title_en_xiaoer_tl { auto }
\tl_new:N \g_hit_titlepage_title_en_xiaoer_tl
\tl_gset:Nn \g_hit_titlepage_title_en_xiaoer_tl { true }
\tl_new:N \g__hit_cover_degree_type_tl
\tl_gset:Nn \g__hit_cover_degree_type_tl { auto }
%    \end{macrocode}
%
% 英文题目的字号，封面与内封各一个同名键。规范给的是二号，题目太长时可用小二，
% 没有第三种字号可选，所以是三值而不是给字号名：
% \begin{itemize}
% \item \texttt{true} 排小二
% \item \texttt{false} 排二号，不降
% \item \texttt{auto} 由类定——封面走让步阶梯（先二号，排不下降小二），内封按
%       题目宽度判
% \end{itemize}
%
% 默认封面 \texttt{auto}、内封 \texttt{false}，与改动前的形态一致。
%    \begin{macrocode}
\keys_define:nn { hit / titlepage }
  {
    title-en-xiaoer .code:n =
      { \__hit_tristate_gset:Nn \g_hit_titlepage_title_en_xiaoer_tl {#1} } ,
    title-en-xiaoer .default:n = true ,
    unknown .code:n =
      { \msg_warning:nnV { hithesis } { unknown-key } \l_keys_key_str } ,
  }
\keys_define:nn { hit / cover }
  {
    title-en-xiaoer .code:n =
      { \__hit_tristate_gset:Nn \g_hit_cover_title_en_xiaoer_tl {#1} } ,
    title-en-xiaoer .default:n = true ,
    degree-type .code:n =
      { \__hit_tristate_set:Nn \g__hit_cover_degree_type_tl {#1} } ,
    degree-type .default:n = true ,
    unknown .code:n =
      { \msg_warning:nnV { hithesis } { unknown-key } \l_keys_key_str } ,
  }
\keys_define:nn { hit / layout }
  {
    linebreak .code:n = { \keys_set:nn { docxlike } { linebreak = {#1} } } ,
    compatibility-mode .code:n =
      { \keys_set:nn { docxlike } { compatibility-mode = {#1} } } ,
    hyphenation .code:n = { \keys_set:nn { docxlike } { hyphenation = {#1} } } ,
  }
\keys_define:nn { hit }
  {
    cover .code:n = { \hit@keys@set:nn { hit / cover } {#1} } ,
    titlepage .code:n = { \hit@keys@set:nn { hit / titlepage } {#1} } ,
    info  .code:n = { \hit@keys@set:nn { hit / info } {#1} } ,
    style .code:n = { \hit@keys@set:nn { hit / style } {#1} } ,
    layout .code:n = { \hit@layout@set:n {#1} } ,
    styles .code:n =
      {
        \hit@styles@set:n {#1}
        % 页眉行盒按 Header 样式算（见 hit-docxlike.dtx），导言区里改了样式就重算一次
        % 版面；正文里不能再动版面，也不需要
        \hit@layout@refresh:
      } ,
    toc   .code:n =
      {
        \str_case:nnF {#1}
          {
            { true }  { \tl_gset:cn { g__hit_structure_table-of-contents_tl } { true } }
            { false } { \tl_gset:cn { g__hit_structure_table-of-contents_tl } { false } }
          }
          { \hit@keys@set:nn { hit / toc } {#1} }
      } ,
    toc .default:n = true ,
    bib   .code:n =
      { \hit@keys@set:nn { hit / bib } {#1} } ,
    term .code:n = { \hit@keys@set:nn { hit / term } {#1} } ,
    const .code:n =
      {
        \__hit_warn_once:nnn { const } { renamed-family } { term }
        \hit@keys@set:nn { hit / term } {#1}
      } ,
    math  .code:n = { \hit@keys@set:nn { hit / math } {#1} } ,
    struct .code:n =
      { \__hit_structure_reset_written: \hit@keys@set:nn { hit / struct } {#1} } ,
    defense .code:n = { \hit@keys@set:nn { hit / defense } {#1} } ,
    abbr .code:n = { \hit@keys@set:nn { hit / abbr } {#1} } ,
    thm .code:n = { \hit@theorem@style {#1} } ,
    unknown .code:n = { \__hit_legacy_key:n {#1} } ,
  }
\RenewDocumentCommand \hitsetup { O{} +m }
  {
    \hit@if@condition:nT {#1}
      { \hit@keys@set:nn { hit } {#2} \hit@bib@backend@apply: }
  }
%    \end{macrocode}
%
% \changes{v3.1d}{2025/03/03}{将 \cs{firstpagecdate} 的设置藏起来}
% 首页那行日期没单独写就跟着封面日期走。判断看的是存原值的
% \cs{hit@cover@date@zh@raw}——\cs{hit@cover@date@zh} 现在是个格式化取值的宏，
% 本身永远非空。直接把原值抄过去，两个字段的精度声明一样，排出来也一样。
%    \begin{macrocode}
\AfterPreamble
  {
    \tl_if_empty:cT { hit@info@cover@date@zh@raw }
      { \cs_gset_eq:cc { hit@info@cover@date@zh@raw } { hit@info@date@zh@raw } }
  }
%</shared-keylist>
%    \end{macrocode}
%
% \section{开题中期报告的键清单（keylist）}
%
% 开题中期报告用到的元信息字段：字段名、旧名映射、\cs{hitsetup} 各族的路由，
% 以及摘要环境。都是清单，不含机制。
%
% ^^A ======================================================================
% ^^A Module report-keylist: 元信息字段（\hitsetup 接受的 key）（hit-report）
% ^^A ======================================================================
%    \begin{macrocode}
%<*shared-keylist>
%% 中文封面

\hit@define@term{report-type-zh}
\hit@define@term{class-id}
\hit@define@date@term{enroll-date-zh}{zh}
%    \end{macrocode}
%
% 中文封面日期的精度在这里重设一遍。共享层那份按终稿的规矩写：本科到年月日、
%    \begin{macrocode}
\hit@define@date@term{date-zh}{zh}

%    \end{macrocode}
%
% 关键词怎么拆、\cs{hitsetup} 的占位定义都见 \file{src/engine/hit-key-engine.dtx}。
%
% 入口换成 l3keys，根族路由到子族，子族一律把键转回
% \texttt{hit} 族。转发时值要重新加花括号，否则含逗号的值会被二次拆分成多个键。
%
% 元信息字段与可本地化字符串。做法与 \file{hit-thesis.cls} 相同：
% 前者的处理器直接调用同名命令，后者的键名由宏名去掉 \texttt{hit@} 前缀、
% 其余 \texttt{@} 换成连字符得来。
%    \begin{macrocode}
\hit@declare@info@field
  {
    class-id , enroll-date-zh , report-type-zh
  }
\hit@declare@constant
  {
    header-report , header-report-front , cover-school-bottom-mark-shenzhen-master ,
    cover-class-bachelor , cover-check-report-title ,
    cover-university-shenzhen-bachelor-zh , cover-university-shenzhen-bachelor-en ,
    cover-form-line-shenzhen-bachelor-zh , cover-form-line-shenzhen-bachelor-en ,
    cover-stage-line-shenzhen-bachelor-zh , cover-stage-line-shenzhen-bachelor-en ,
    cover-title-label-shenzhen-bachelor-zh , cover-title-label-shenzhen-bachelor-en ,
    cover-author-label-shenzhen-bachelor-zh , cover-author-label-shenzhen-bachelor-en ,
    cover-student-id-label-shenzhen-bachelor-zh , cover-student-id-label-shenzhen-bachelor-en ,
    cover-affiliation-label-shenzhen-bachelor-zh , cover-affiliation-label-shenzhen-bachelor-en ,
    cover-speciality-label-shenzhen-bachelor-zh , cover-speciality-label-shenzhen-bachelor-en ,
    cover-supervisor-label-shenzhen-bachelor-zh , cover-supervisor-label-shenzhen-bachelor-en ,
    cover-date-label-shenzhen-bachelor-zh , cover-date-label-shenzhen-bachelor-en ,
    cover-university-font , cover-document-type-font , cover-report-type-line ,
    cover-degree-prefix
  }
%    \end{macrocode}
%
% 参考文献的三个选项，改用连字符命名。
%
% 键名、设值器与守卫跟终稿同一套，见那边的说明。报告没有
% \option{compress-min} 与 \option{split-entry}，跟统一前一样。
%    \begin{macrocode}
%    \end{macrocode}
%
% 共享层的 \option{toc} 族三个键，报告只做得了两个。一份报告只排一份目录，
% 语言跟着类选项 \option{lang} 走，没有「中文目录后面再跟一份英文的」这回事。
% 这里把 \option{english} 覆盖成一条提示，比让它静静地不起作用强。
%    \begin{macrocode}
%    \end{macrocode}
%
% 兼容期的处理，与 \file{hit-thesis.cls} 一致：旧名照常生效，只在日志里提示。
%    \begin{macrocode}
%    \end{macrocode}
%
% 旧名到新名的对照。旧名在 v3.1e 就是公开接口，写在用户的 \file{cover.tex} 里，
% 所以保留到 v4。
%    \begin{macrocode}
\hit@define@legacy@terms
  {
    ctype=report-type-zh ,
    csubject=major-zh ,
    cauthor=author-zh ,
    cstudentid=student-id ,
    cclassid=class-id ,
    caffil=affiliation-zh ,
    csupervisor=supervisor-zh ,
    cdate=date-zh ,
    cenrolldate=enroll-date-zh
  }
\hit@define@legacy@title{ctitlecover}{zh}{first}
\hit@define@legacy@title{ctitleone}{zh}{first}
\hit@define@legacy@title{ctitletwo}{zh}{second}
%    \end{macrocode}
%
% 根族路由。\option{toc} 既是族名又是老的布尔选项名，所以它单独处理：
% 值是 \texttt{true}/\texttt{false} 时当老写法，其余当子族的键值表。
%    \begin{macrocode}
%</shared-keylist>
%    \end{macrocode}
%
% \Finale
